\documentclass[11pt]{article}

\usepackage{fullpage}
\usepackage{amsmath}
\usepackage{amssymb}
\usepackage{amsfonts}
\usepackage{amsthm}
\usepackage{mathtools}
\usepackage{enumitem}
\usepackage[colorlinks,linkcolor=magenta,citecolor=blue]{hyperref}

\newcommand{\Rbar}{\bar R}
\newcommand{\Om}{\Omega}
\newcommand{\E}{\mathbb E}
\newcommand{\Var}{\operatorname{Var}}
\newcommand{\tr}{\operatorname{tr}}
\newcommand{\diag}{\operatorname{diag}}
\newcommand{\Unif}{\operatorname{Unif}}
\newcommand{\IG}{\operatorname{IG}}
\newcommand{\Gam}{\operatorname{Gamma}}
\newcommand{\Normal}{\mathcal N}

\newtheorem{lemma}{Lemma}

\title{Susie-relax with the ``irrelevance of the null SNPs'' property}
\author{}
\date{}

\begin{document}
\maketitle

\section{SER-relax-prime}
Now we consider another approximation to truly achieve the ``irrelevance of the null SNPs'' property. It leads to another model that I call `SER-relax-prime'.

Recall the SER-relax (conditional on $\gamma = 1$):

\begin{equation}\label{eq:SER-gamma-1}
x \mid \beta,\,\gamma=1 \;\sim\;
\mathcal{N}\!\left(
  \beta\begin{pmatrix}1\\\mu_1\end{pmatrix},\;
  \frac{1}{N}\begin{pmatrix}1 & \mu_1^T \\ \mu_1 & \bar{R}_{-1,-1}+N\beta^2\tau^2 S_1\end{pmatrix}
\right)
\end{equation}
\paragraph{Matthew + Gemini suggest} omitting $\beta^2$ (replacing $\beta^2$ by $1$) from the variance in the formula above. Similar to computation in PIP from the previous file, we have

\begin{align}\label{eq:conditional-prob-x-given-beta-gamma-prime}
p(x \mid \beta,\,\gamma=j)
&= \frac{N^{J/2}}{(2\pi)^{J/2}\,|\bar{R}|^{1/2}\,(1+N\tau^2)^{(J-1)/2}}
\exp\!\left(-\frac{N}{2(1+N\tau^2)}q\right) \nonumber \\
&\quad\times
\exp\!\left(-\frac{N^2\tau^2}{2(1+N\tau^2)}x_j^2\right)
\exp\!\left(N\beta x_j - \frac{N}{2}\beta^2\right).
\end{align}
Integrating out $\beta \sim \mathcal{N}(0, \sigma_{0}^2)$, we have
\begin{align}\label{eq:conditional-prob-x-given-gamma-prime}
p(x \mid \gamma=j)
&= \frac{N^{J/2}}{(2\pi)^{J/2}\,|\bar{R}|^{1/2}\,(1+N\tau^2)^{(J-1)/2}}
\exp\!\left(-\frac{N}{2(1+N\tau^2)}q\right) \nonumber \\
&\quad\times \sqrt{\dfrac{1}{1 + N\sigma_{0}^2}} \times \exp\left(\left(\dfrac{N \sigma_0^2}{1 + N \sigma_0^2} - \dfrac{N\tau^2}{1 + N \tau^2} \right) \dfrac{N x_j^2}{2} \right).
\end{align}
Leading to
\begin{align}\label{eq:PIP-SER-prime-unit}
p(\gamma=j | x) 
\propto \exp\left(\left(\dfrac{N \sigma_0^2}{1 + N \sigma_0^2} - \dfrac{N\tau^2}{1 + N \tau^2} \right) \dfrac{N x_j^2}{2} \right).
\end{align}
So we truly have the desired property. However, something is not right about the formula above: The scalar in the exponential can turn into negative if $\sigma_0^2 < \tau^2$ (so that a smaller $Z$-score leads to a bigger PIP?). It is because we are replacing $\beta^2$ by $1$ in the variance of the likelihood of $p(x | \beta, \gamma = j)$. Maybe the approximation $\beta^2 \approx \sigma_0^2$ makes more sense.

\paragraph{SER-relax-prime.} This model approximates $\beta^2$ in the variance of~\eqref{eq:SER-gamma-1} by $\sigma_{0}^2$. It leads to 
\begin{equation}\label{eq:SER-prime-gamma-1}
x \mid \beta,\,\gamma=1 \;\sim\;
\mathcal{N}\!\left(
  \beta\begin{pmatrix}1\\\mu_1\end{pmatrix},\;
  \frac{1}{N}\begin{pmatrix}1 & \mu_1^T \\ \mu_1 & \bar{R}_{-1,-1}+N\sigma_0^2\tau^2 S_1\end{pmatrix}
\right),
\end{equation}
and similarly for other $\gamma = j$ (recall $\mu_1 = \bar{R}_{-1, 1}$). The PIP in this model is:
\begin{align}\label{eq:PIP-SER-prime}
p(\gamma=j | x) 
\propto \exp\left(\left(\dfrac{N \sigma_0^2}{1 + N \sigma_0^2} - \dfrac{N \sigma_0^2 \tau^2}{1 + N \sigma_0^2 \tau^2} \right) \dfrac{N x_j^2}{2} \right).
\end{align}
Note that $\tau^2 = 1 / (\nu_0 + 3) \leq 1/ 3 < 1$, so that $N \sigma_0^2 \tau^2 < N \sigma_0^2$, implying
$$ \dfrac{N \sigma_0^2}{1 + N \sigma_0^2} - \dfrac{N \sigma_0^2 \tau^2}{1 + N \sigma_0^2 \tau^2} > 0.$$
This quantity tends to $\dfrac{N \sigma_0^2}{1 + N \sigma_0^2}$ as $\tau^2 \to 0$ (recovering susie-rss). One interpretation of~\eqref{eq:PIP-SER-prime} is that PIP might be more a little more diffuse when we are more uncertain about the in-sample $R$ (which can be a bit funny when we think about it, because if $L = 1$ the LD does not matter and the PIP ideally should be similar to susie-rss. But I guess nothing is perfect, and this really shows how Susie-relax does ``relax'' the signal.) 

However, one can show that this PIP retains some important properties:

\begin{enumerate}
\item PIP of SNP j is exponentially proportional to the Z-score square
\item if two SNPs have similar Z-score, they will have the same PIP
\item To compare this PIP with susie-rss, write
\[
t=N\sigma_0^2,\qquad s=\tau^2.
\]
The susie-rss coefficient in the exponent is
\[
a_{\mathrm{rss}}=\frac{t}{1+t},
\]
whereas the SER-relax-prime coefficient is
\[
a_{\mathrm{prime}}
=
\frac{t}{1+t}-\frac{ts}{1+ts}.
\]
Thus
\[
0\le a_{\mathrm{rss}}-a_{\mathrm{prime}}
=
\frac{ts}{1+ts}
\le ts.
\]
If $\nu_0\ge 10$, then $s=\tau^2=1/(\nu_0+3)\le 1/13$, so
\[
a_{\mathrm{rss}}-a_{\mathrm{prime}}
\le
\frac{t}{13+t}
\le
\frac{t}{13}.
\]
Equivalently, for any two SNPs $j$ and $k$,
\[
\left|
\log\frac{p_{\mathrm{prime}}(\gamma=j\mid x)}
{p_{\mathrm{prime}}(\gamma=k\mid x)}
-
\log\frac{p_{\mathrm{rss}}(\gamma=j\mid x)}
{p_{\mathrm{rss}}(\gamma=k\mid x)}
\right|
=
\frac{N}{2}\frac{ts}{1+ts}|x_j^2-x_k^2|
\le
\frac{N}{2}\frac{t}{13+t}|x_j^2-x_k^2|.
\]
Hence, when $t=N\sigma_0^2$ is not too large, $\nu_0=10$ already gives only a
small flattening of the susie-rss log-odds. The approximation becomes exact as
$\nu_0\to\infty$.
\end{enumerate}

\paragraph{Posterior distribution of $\beta$ under SER-relax-prime.}
Under~\eqref{eq:SER-prime-gamma-1}, the covariance no longer depends on $\beta$.
Therefore, conditional on $\gamma=j$, the likelihood as a function of $\beta$ is
\[
p(x\mid \beta,\gamma=j)
\propto
\exp\!\left(N\beta x_j-\frac{N}{2}\beta^2\right),
\]
where all remaining terms are constant in $\beta$. Combining this likelihood
with the prior $\beta\sim\mathcal N(0,\sigma_0^2)$ gives
\begin{align}
p(\beta\mid x,\gamma=j)
&=
\mathcal N(m_j,v),\\
v&=\left(N+\sigma_0^{-2}\right)^{-1},\\
m_j&=vNx_j.
\end{align}
Thus the conditional posterior variance is shared across coordinates, and the
posterior mean depends on the data only through the active marginal statistic
$x_j$. In particular, $p(\beta\mid x,\gamma=j)$ does not depend on
$q=x^T\bar R^{-1}x$.

\section{susie-relax-prime}
How to extend~\eqref{eq:SER-prime-gamma-1} to the sum of a single-effect type of model? Ideally, each of the single-effect models should have the same noise model (like $\mathcal{N}(0, \bar R / N)$) to enable IBSS-like inference. I think it is cleanest to introduce a latent variable $\eta_{\ell}$ and an
effect contribution vector $\theta_\ell$ for each single-effect component.
For $\ell=1,\ldots,L$,
\[
\gamma_\ell\sim\Unif\{1,\ldots,J\},\qquad
\beta_\ell\sim\mathcal N(0,\sigma_{0\ell}^2).
\]
Conditional on $\gamma_\ell=j$, define
\[
\theta_\ell=\beta_\ell\bar R_{\cdot j}+\eta_{\ell j},
\]
where $\eta_{\ell j,j}=0$ and
\[
\eta_{\ell j,-j}\sim\mathcal N(0,\sigma_{0\ell}^2\tau^2S_j).
\]
The observation model is
\[
x\mid \theta_1,\ldots,\theta_L
\sim
\mathcal N\!\left(
\sum_{\ell=1}^L\theta_\ell,\frac{\bar R}{N}
\right).
\]
This is a proper generative model. If we integrate out $\eta_{\ell j}$ for a
single effect, we recover the SER-relax-prime covariance in
\eqref{eq:SER-prime-gamma-1}, with $\sigma_0^2$ replaced by
$\sigma_{0\ell}^2$.

\paragraph{Variational family and ELBO.}
Use the mean-field approximation
\[
q=
\prod_{\ell=1}^L
q_\ell(\gamma_\ell,\beta_\ell,\eta_\ell),
\]
with
\[
q_\ell(\gamma_\ell=j,\beta_\ell,\eta_{\ell j})
=
\alpha_{\ell j}
q_{\ell j}(\beta_\ell)
q_{\ell j}(\eta_{\ell j}).
\]
The ELBO is the standard variational objective
\begin{align}
\mathcal L(q)
&=
E_q\log p(x\mid \theta_1,\ldots,\theta_L)
+\sum_{\ell=1}^L E_q\log p(\gamma_\ell,\beta_\ell,\eta_\ell\mid\gamma_\ell)
\nonumber\\
&\quad
-\sum_{\ell=1}^L E_q\log q_\ell(\gamma_\ell,\beta_\ell,\eta_\ell).
\end{align}
Holding all factors except $q_\ell$ fixed, the optimal coordinate update is
\[
\log q_\ell^*(\gamma_\ell,\beta_\ell,\eta_\ell)
=
E_{-\ell}\log p(x,\{\gamma_k,\beta_k,\eta_k\}_{k=1}^L)+\text{const}.
\]
Let
\[
\bar\theta_k=E_q(\theta_k),\qquad
r_\ell=x-\sum_{k\ne\ell}\bar\theta_k.
\]
Then the $q_\ell$ update is exactly the single-effect SER-relax-prime posterior
computed with $x$ replaced by $r_\ell$.

\paragraph{Coordinate posterior.}
For each effect define
\[
\kappa_\ell=\frac{N\sigma_{0\ell}^2}{1+N\sigma_{0\ell}^2},
\qquad
\rho_\ell=\frac{N\sigma_{0\ell}^2\tau^2}{1+N\sigma_{0\ell}^2\tau^2}.
\]
The coordinate weights are
\begin{align}
\alpha_{\ell j}
&=
\frac{\exp\!\left\{\dfrac{N}{2}(\kappa_\ell-\rho_\ell)r_{\ell j}^2\right\}}
{\sum_{k=1}^J
\exp\!\left\{\dfrac{N}{2}(\kappa_\ell-\rho_\ell)r_{\ell k}^2\right\}}.
\end{align}
The common factor involving
$q_\ell^{\mathrm{quad}}=r_\ell^T\bar R^{-1}r_\ell$ cancels in this
normalization. Conditional on $\gamma_\ell=j$,
\begin{align}
q_{\ell j}(\beta_\ell)
&=
\mathcal N(m_{\ell j},v_\ell),\\
v_\ell&=\left(N+\sigma_{0\ell}^{-2}\right)^{-1},\\
m_{\ell j}&=v_\ell N r_{\ell j}
=\kappa_\ell r_{\ell j}.
\end{align}
The relaxed part also has a Gaussian posterior. Its nonzero block is
\begin{align}
q_{\ell j}(\eta_{\ell j,-j})
&=
\mathcal N(\bar\eta_{\ell j,-j},W_{\ell j}),\\
\bar\eta_{\ell j,-j}
&=
\rho_\ell(r_{\ell,-j}-\mu_jr_{\ell j}),\\
W_{\ell j}
&=
\frac{\sigma_{0\ell}^2\tau^2}{1+N\sigma_{0\ell}^2\tau^2}S_j.
\end{align}
Equivalently, as a $J$-vector,
\[
\bar\eta_{\ell j}
=
\rho_\ell(r_\ell-r_{\ell j}\bar R_{\cdot j}).
\]

\paragraph{Posterior mean contribution.}
The posterior mean contribution of effect $\ell$ must include both the
posterior mean of $\beta_\ell\bar R_{\cdot j}$ and the posterior mean of the
relaxed component $\eta_{\ell j}$:
\begin{align}
\bar\theta_\ell
&=
\sum_{j=1}^J
\alpha_{\ell j}
\left[
m_{\ell j}\bar R_{\cdot j}
+\rho_\ell(r_\ell-r_{\ell j}\bar R_{\cdot j})
\right] \nonumber\\
&=
\bar R(\alpha_\ell\odot m_\ell)
+\rho_\ell r_\ell
-\rho_\ell\bar R(\alpha_\ell\odot r_\ell).
\end{align}
This correction is what makes the residual update consistent with the
variational posterior under the latent contribution model.

\paragraph{ELBO computation.}
Let $\Omega=\bar R^{-1}$ and
\[
C_x=-\frac{J}{2}\log(2\pi)-\frac12\log\left|\frac{\bar R}{N}\right|.
\]
For each $(\ell,j)$ define
\[
c_\ell
=
\frac{\sigma_{0\ell}^2\tau^2}
{1+N\sigma_{0\ell}^2\tau^2},
\qquad
C_{\ell j}
=
(r_{\ell,-j}-\mu_jr_{\ell j})^TS_j^{-1}
(r_{\ell,-j}-\mu_jr_{\ell j}).
\]
Using the block inverse identity,
\[
C_{\ell j}=r_\ell^T\Omega r_\ell-r_{\ell j}^2,
\]
so all $C_{\ell j}$ can be computed from the scalar
$r_\ell^T\Omega r_\ell$ and the coordinate-wise residuals. The posterior
second moments are
\begin{align}
B_{\ell j}
&=
E_q(\beta_\ell^2\mid\gamma_\ell=j)
=m_{\ell j}^2+v_\ell,\\
D_{\ell j}
&=
E_q(\eta_{\ell j,-j}^TS_j^{-1}\eta_{\ell j,-j})
=
\rho_\ell^2C_{\ell j}+(J-1)c_\ell.
\end{align}
Because $\bar R_{\cdot j}^T\Omega=e_j^T$ and $\eta_{\ell j,j}=0$, the cross
term between $\beta_\ell\bar R_{\cdot j}$ and $\eta_{\ell j}$ is zero. Hence
\[
E_q(\theta_\ell^T\Omega\theta_\ell)
=
\sum_{j=1}^J\alpha_{\ell j}(B_{\ell j}+D_{\ell j}).
\]

Let $\bar\theta=\sum_{\ell=1}^L\bar\theta_\ell$. A fast way to compute the
expected quadratic loss is
\[
E_q\left[
(x-\sum_{\ell=1}^L\theta_\ell)^T
\Omega
(x-\sum_{\ell=1}^L\theta_\ell)
\right]
=
(x-\bar\theta)^T\Omega(x-\bar\theta)
+\sum_{\ell=1}^L V_\ell,
\]
where
\[
V_\ell
=
\sum_{j=1}^J\alpha_{\ell j}(B_{\ell j}+D_{\ell j})
-\bar\theta_\ell^T\Omega\bar\theta_\ell.
\]
Thus the expected log-likelihood is
\[
E_q\log p(x\mid\theta_1,\ldots,\theta_L)
=
C_x-\frac{N}{2}
\left[
(x-\bar\theta)^T\Omega(x-\bar\theta)
+\sum_{\ell=1}^L V_\ell
\right].
\]

The remaining terms are most compactly written as KL divergences. Define
\begin{align}
\mathrm{KL}^{\beta}_{\ell j}\label{eq:KL_beta}
&=
\frac12
\left[
\log\frac{\sigma_{0\ell}^2}{v_\ell}
+\frac{B_{\ell j}}{\sigma_{0\ell}^2}
-1
\right],
\end{align}
and
\begin{align}\label{eq:KL_eta}
\mathrm{KL}^{\eta}_{\ell j}
&=
\frac12
\left[
(J-1)\log\frac{\sigma_{0\ell}^2\tau^2}{c_\ell}
+\frac{D_{\ell j}}{\sigma_{0\ell}^2\tau^2}
-(J-1)
\right] \nonumber\\
&=
\frac12
\left[
(J-1)\log(1+N\sigma_{0\ell}^2\tau^2)
+\frac{D_{\ell j}}{\sigma_{0\ell}^2\tau^2}
-(J-1)
\right].
\end{align}
The determinants $|S_j|$ cancel in $\mathrm{KL}^{\eta}_{\ell j}$ because
$W_{\ell j}=c_\ell S_j$. The full ELBO is therefore
\begin{align}\label{eq:full-ELBO}
\mathcal L(q)
&=
C_x-\frac{N}{2}
\left[
(x-\bar\theta)^T\Omega(x-\bar\theta)
+\sum_{\ell=1}^L V_\ell
\right] \nonumber\\
&\quad
-\sum_{\ell=1}^L\sum_{j=1}^J
\alpha_{\ell j}
\left[
\log(J\alpha_{\ell j})
+\mathrm{KL}^{\beta}_{\ell j}
+\mathrm{KL}^{\eta}_{\ell j}
\right].
\end{align}
This is the expression to use in code. It only requires:
\begin{enumerate}[leftmargin=*]
\item the fitted mean $\bar\theta$ and one quadratic
$(x-\bar\theta)^T\Omega(x-\bar\theta)$;
\item for each effect, $\bar\theta_\ell^T\Omega\bar\theta_\ell$;
\item for each effect, $r_\ell^T\Omega r_\ell$, which gives all
$C_{\ell j}=r_\ell^T\Omega r_\ell-r_{\ell j}^2$;
\item coordinate-wise scalar operations for $B_{\ell j}$, $D_{\ell j}$, and the
KL terms.
\end{enumerate}
No determinant of $S_j$ is needed, and no dense $W_{\ell j}$ matrix has to be
formed.

\paragraph{Variational empirical Bayes updates.}
Holding the variational distribution fixed, maximize the ELBO over
$\sigma_{0\ell}^2$ and $\tau^2$. Define the effect-level averages
\[
B_\ell=\sum_{j=1}^J\alpha_{\ell j}B_{\ell j},
\qquad
D_\ell=\sum_{j=1}^J\alpha_{\ell j}D_{\ell j}.
\]
The part of the ELBO that depends on $\sigma_{0\ell}^2$ is
\[
-\frac{J}{2}\log\sigma_{0\ell}^2
-\frac{1}{2\sigma_{0\ell}^2}
\left(B_\ell+\frac{D_\ell}{\tau^2}\right),
\]
so the closed-form VEB update is
\[
\hat\sigma_{0\ell}^2
=
\frac{1}{J}
\left(B_\ell+\frac{D_\ell}{\tau^2}\right).
\]
This is the average posterior second moment of the latent effect: one degree of
freedom comes from $\beta_\ell$, and $J-1$ degrees of freedom come from the
relaxed LD-column perturbation after rescaling by $\tau^2$.

Holding $\{\sigma_{0\ell}^2\}_{\ell=1}^L$ fixed, the $\tau^2$-dependent part of
the ELBO is
\[
-\frac{L(J-1)}{2}\log\tau^2
-\frac{1}{2\tau^2}
\sum_{\ell=1}^L\frac{D_\ell}{\sigma_{0\ell}^2},
\]
so
\[
\hat\tau^2
=
\frac{1}{L(J-1)}
\sum_{\ell=1}^L\frac{D_\ell}{\sigma_{0\ell}^2}.
\]
This update estimates the average standardized size of the relaxed LD-column
perturbations. A small value means the posterior does not need much LD
relaxation and the method approaches susie-rss; a large value means the data
are being explained partly through LD-column uncertainty. In practice these EB
updates can be alternated with the variational updates and bounded to avoid
degenerate values.

\paragraph{IBSS-like coordinate ascent.}
The resulting algorithm is coordinate ascent on $\mathcal L(q)$:
\begin{enumerate}[leftmargin=*]
\item Initialize $\alpha_{\ell j}$, $m_{\ell j}$, $v_\ell$,
$\bar\eta_{\ell j}$, and $\bar\theta_\ell$ for $\ell=1,\ldots,L$.
\item For each effect $\ell$, form
\[
r_\ell=x-\sum_{k\ne\ell}\bar\theta_k.
\]
\item Compute $v_\ell=(N+\sigma_{0\ell}^{-2})^{-1}$,
$m_{\ell j}=v_\ell Nr_{\ell j}$, and
$\bar\eta_{\ell j}=\rho_\ell(r_\ell-r_{\ell j}\bar R_{\cdot j})$ for all $j$.
\item Update the coordinate weights by
\[
\alpha_{\ell j}
\propto
\exp\!\left\{
\frac{N}{2}(\kappa_\ell-\rho_\ell)r_{\ell j}^2
\right\}.
\]
\item Recompute
\[
\bar\theta_\ell
=
\bar R(\alpha_\ell\odot m_\ell)
+\rho_\ell r_\ell
-\rho_\ell\bar R(\alpha_\ell\odot r_\ell).
\]
\item Optionally update $\sigma_{0\ell}^2$ and $\tau^2$ using the VEB formulas,
then refresh $\kappa_\ell$, $\rho_\ell$, $v_\ell$, and $W_{\ell j}$.
\item Sweep over $\ell=1,\ldots,L$ until the ELBO or fitted mean
$\sum_{\ell=1}^L\bar\theta_\ell$ stabilizes.
\end{enumerate}

\paragraph{Comment on the robustness.} We knew from the experiments that under LD mismatch, Susie-RSS's signal $R \bar{b}$ can't fit exactly to the data $x$ (previously denoted by $v$) and therefore needs more $L$. This model allow for some deviation from $R \bar{b}$ by additionally using $(\eta_{\ell})_{\ell = 1}^{L}$. Looking at the final ELBO~\eqref{eq:full-ELBO}, the choice of $L$ now is a battle between $\tau^2$ and $\sigma_{0\ell}^2$ in equations~\eqref{eq:KL_beta} and~\eqref{eq:KL_eta}: If introduce a little more $\tau^2$ can save us the cost of one more $\sigma_{0\ell}^2$, the model will rather set $\tau^2$ to slightly larger (allowing more deviation) and set that extra $\sigma_{0\ell}^2$ to 0, leading to more conservative behavior.
\end{document}
